\section{Introduction}
\label{sec:intro}

\IEEEPARstart{R}{iver} and reservoir monitoring with in-situ sensor nodes has to work where field visits are expensive and power is scarce~\cite{adumanu2017survey,adumanu2020smart}. A typical node samples temperature (T), specific conductance (EC), dissolved oxygen (DO) and pH every 5--30~min and sends the readings over a cellular or low-power wide-area network (LPWAN). Such nodes lose data. The InWaterSense deployment on the Sitnica river in Kosovo missed \inwsLoss{} of its expected readings~\cite{ahmedi2022dataset}; the CSIRO water-quality network in Australia averaged about 64\% delivery because the gateway's connection to the server ``was down occasionally''~\cite{ledinh2007design}; and deployment reports describe gateways that lose power every night~\cite{langendoen2006murphy}, sinks that stop reconnecting until rebooted~\cite{ingelrest2010sensorscope,brzozowski2018lessons}, and outages of days to weeks~\cite{puluckul2025subarctic,zennaro2025kyrgyzstan}.

A node can \emph{retransmit} a lost reading until it is acknowledged (automatic repeat request, ARQ), \emph{buffer} it and backfill it later, or leave the server to \emph{reconstruct} it from the site's history or its neighbours. Which response is right depends on how data are lost: retransmission helps against a short radio fade but wastes energy against a gateway that is down for a day, and reconstruction can fill ten minutes but not a week. Yet the analyses behind these remedies assume short losses: estimation over lossy links uses independent or short-memory (Markov fading) packet errors~\cite{sinopoli2004kalman,huang2019retransmit,huang2020harq}, and imputation methods are evaluated on random masks or simulated failures of a few hours~\cite{kong2014dataloss,cini2022grin,marisca2022spin,chang2023deepcs}. How long real delivery gaps last, how much of the missing data they hold, and what this implies for the energy and timeliness of each remedy have not been measured for environmental telemetry.

The contribution of this paper is not a new retransmission or prediction primitive; ARQ, buffering, back-off, batching and error-bounded (dual-prediction) reporting are established. It is an evidence-to-design chain for outage-aware, energy-efficient telemetry: real delivery records measure how long losses last and where the missing data are; the records are replayed through a measurement-parameterised LPWAN energy model, with explicit bounds on what they cannot show (why a reading is missing, and when within a slot the network returned), to decide when each remedy pays; and the resulting reporting design is fixed on one year of river data and tested on a held-out year. Four research questions structure the paper:
\begin{enumerate}
\item[RQ1] In the records studied, how long are the delivery gaps, how much of the missing data do long gaps hold, and is a geometric (memoryless) gap-length model adequate?
\item[RQ2] Under these gaps, what do retransmission, buffering with back-off and batching recover, at what energy and delay, on LoRaWAN and NB-IoT?
\item[RQ3] Can reconstruction from a site's own history or its neighbours replace delivery during a gap?
\item[RQ4] Can one reporting scheme raise alarms as fast as per-reading reporting and keep the server's real-time picture within a stated tolerance as often, at substantially less energy than per-reading retransmission, on data not used to design it?
\end{enumerate}

The contributions follow the chain.
\begin{itemize}
\item \emph{Empirical loss structure} (Sec.~\ref{sec:outages}). In four records, two water-quality deployments and two LoRaWAN networks, losses separate into short radio losses and rare long delivery gaps, which hold most of the missing readings in the water-quality records. A geometric gap-length model is strongly unsupported relative to the alternatives and fails worst in the long-gap tail. Where the records allow it, the cause of long gaps is traced.
\item \emph{Consequence analysis with explicit bounds} (Secs.~\ref{sec:handling} and~\ref{sec:reconstruct}). ARQ recovers under a fifth of the missing readings even under optimistic retry timing, at \arqRatioMin--\arqRatioMax{} the energy of unconfirmed reporting on LoRaWAN, whereas buffering with back-off recovers every generated reading at an energy close to that of recovering nothing. The results are bracketed for the unknown cause of gaps, retry timing, storage, the duty cycle, model uncertainty and many nodes recovering at once. The evaluated reconstruction estimators, including two deep imputers, keep few day-long gaps within tolerance.
\item \emph{An outage-aware reporting design, confirmed on held-out data} (Sec.~\ref{sec:design}). Bounded event-driven batching combines hourly batches, immediate alarm uploads, a bounded-error trigger that mirrors the server's estimate, and buffering with back-off. It was developed on \designSites{} USGS river sites in 2023 and tested, with its rules frozen, on \confSites{} sites in 2024.
\end{itemize}

\emph{Scope.} The two water-quality records cannot fully separate a network outage from a node that stopped producing readings, so the paper speaks of delivery gaps and treats buffering's recovery as an upper bound. Energy comes from a model parameterised by published measurements, not from hardware built for this study, and the claims are ratios between strategies on the same radio. A supplement gives the details behind each section.

Sec.~\ref{sec:related} reviews related work and Sec.~\ref{sec:data} the data and models; Secs.~\ref{sec:outages}--\ref{sec:design} answer RQ1--RQ4, and Sec.~\ref{sec:discussion} discusses guidance and limitations.
